% =============================================================
% 06_network.tex — ネットワークとデータベース
% =============================================================
\chapter{ネットワークとデータベース}

\chaptermeta{%
  \item TCP/IPの4層モデルを説明できる
  \item 主要プロトコルの役割を即答できる
  \item SQLの基本構文（SELECT/WHERE/ORDER BY）が読める
  \item 正規化の目的と第1〜第3正規形を区別できる
}{90分}{3}{第5章}

テクノロジ系の中盤は「つながる仕組み」と「データの管理」です。
TCP/IPの階層は\textbf{プロトコルの分類問題}で頻出し、
SQLは\textbf{文法を読めるかどうか}がそのまま得点に直結します。

% =============================================================
\section{TCP/IPの4層モデル}
% =============================================================

\begin{teigi}[TCP/IP 4層]
\begin{enumerate}
  \item \termtag{アプリケーション層}：HTTP, HTTPS, SMTP, POP3, FTP, DNS, DHCP, NTP
  \item \termtag{トランスポート層}：TCP, UDP
  \item \termtag{インターネット層}：IP, ICMP, ARP
  \item \termtag{ネットワークインターフェース層}：Ethernet, Wi-Fi
\end{enumerate}
OSI参照モデル（7層）との対応も問われるが、Iパスでは4層が中心。
\end{teigi}

\begin{hikaku}[主要プロトコルの整理]
\comptable{プロトコル & 層 & 役割}{
HTTP/HTTPS & アプリケーション & Webページの送受信（Sは暗号化あり） \\
SMTP & アプリケーション & メール送信 \\
POP3/IMAP & アプリケーション & メール受信（POP3=DL型、IMAP=サーバ参照型） \\
FTP & アプリケーション & ファイル転送 \\
DNS & アプリケーション & ドメイン名 ↔ IPアドレスの変換 \\
DHCP & アプリケーション & IPアドレスの自動割り当て \\
NTP & アプリケーション & 時刻同期 \\
TCP & トランスポート & 信頼性のある通信（コネクション型） \\
UDP & トランスポート & 高速だが信頼性なし（コネクションレス型）
}
\end{hikaku}

% =============================================================
\section{TCPとUDP}
% =============================================================

\begin{hikaku}[TCP vs UDP]
\comptable{特性 & TCP & UDP}{
通信方式 & コネクション型 & コネクションレス型 \\
信頼性 & 高い（再送制御あり） & 低い（再送なし） \\
速度 & 相対的に遅い & 高速 \\
用途 & Web・メール・ファイル転送 & 動画配信・VoIP・DNS問い合わせ
}
\end{hikaku}

% --- 例題 ---
\begin{reidai}{TCPとUDPの識別}
リアルタイムの動画配信に適したトランスポート層のプロトコルはどれか。
信頼性よりも速度を重視する場合を考えよ。

\sentakushi{FTP}{HTTP}{TCP}{UDP}
\end{reidai}

\begin{shishin}
「リアルタイム」「速度重視」がキーワード。FTPとHTTPはアプリケーション層で、トランスポート層ではない。
TCPは信頼性重視で再送制御がある分、遅延が発生しやすい。
UDPは再送せず高速だが、データが欠落することがある。
動画配信では多少の欠落よりリアルタイム性を優先する。
\end{shishin}

\begin{kaitou}
\textbf{正解：エ（UDP）}

UDPはコネクションレス型で再送制御がなく高速。
リアルタイム配信・VoIP・オンラインゲームなど、遅延より速度を重視する場面で使われる。
\end{kaitou}

% =============================================================
\section{DNS・DHCP}
% =============================================================

\begin{teigi}[DNSとDHCP]
\begin{itemize}
  \item \termtag{DNS}（Domain Name System）：ドメイン名（www.example.com）をIPアドレス（192.168.1.1）に変換する「電話帳」
  \item \termtag{DHCP}（Dynamic Host Configuration Protocol）：ネットワーク接続時にIPアドレスを自動的に割り当てる
\end{itemize}
\end{teigi}

\begin{pitfall}[DNSとDHCPの混同]
\begin{itemize}
  \item DNS → \textbf{名前解決}（ドメイン名 ↔ IPアドレスの変換）
  \item DHCP → \textbf{アドレス配布}（IPアドレスの自動割り当て）
  \item 「DNSは翻訳、DHCPは配布」と覚える
\end{itemize}
\end{pitfall}

% =============================================================
\section{サブネットマスク}
% =============================================================

\begin{teigi}[IPアドレスとサブネットマスク]
\begin{itemize}
  \item \termtag{IPアドレス}：ネットワーク上の機器を識別する番号（IPv4は32ビット）
  \item \termtag{サブネットマスク}：IPアドレスの「ネットワーク部」と「ホスト部」の境界を示す
  \item 例：255.255.255.0 → 上位24ビットがネットワーク部、下位8ビットがホスト部
  \item 同じネットワーク部を持つ機器同士が「同一ネットワーク」
\end{itemize}
\end{teigi}

\begin{pitfall}[サブネットマスクの計算ミス]
サブネットマスク255.255.255.0の場合、ホスト部は8ビット。
使えるアドレス数は $2^8 - 2 = 254$（ネットワークアドレスとブロードキャストを除く）。
「$-2$」を忘れがち。
\end{pitfall}

% =============================================================
\section{RDB（関係データベース）の基礎}
% =============================================================

\begin{teigi}[RDBの基本用語]
\begin{itemize}
  \item \termtag{テーブル（表）}：データを行と列で格納する
  \item \termtag{レコード（行/タプル）}：1件分のデータ
  \item \termtag{フィールド（列/属性）}：データの項目
  \item \termtag{主キー}：各レコードを一意に識別するフィールド（NULLや重複不可）
  \item \termtag{外部キー}：他のテーブルの主キーを参照するフィールド（参照整合性の確保）
\end{itemize}
\end{teigi}

\begin{hikaku}[主キーと外部キー]
\comptable{項目 & 主キー & 外部キー}{
役割 & 自テーブルのレコードを一意に識別 & 他テーブルの主キーを参照 \\
NULL & 不可 & 許可される場合あり \\
重複 & 不可 & 許可される場合あり \\
目的 & レコードの特定 & テーブル間の関連づけ
}
\end{hikaku}

% =============================================================
\section{SQL}
% =============================================================

\begin{teigi}[SQLの基本構文]
\begin{itemize}
  \item \termtag{SELECT}：取得する列を指定
  \item \termtag{FROM}：対象テーブルを指定
  \item \termtag{WHERE}：条件で絞り込み
  \item \termtag{ORDER BY}：並び替え（ASC=昇順、DESC=降順）
  \item \termtag{GROUP BY}：グループ化（集計関数 COUNT, SUM, AVG, MAX, MIN と併用）
  \item \termtag{HAVING}：GROUP BYの結果に対する条件
\end{itemize}
\end{teigi}

\begin{advice}[SQLは「読めればOK」]
Iパスでは SQL を書く問題は出ず、\textbf{SELECT文を読んで結果を選ぶ}形式。
WHERE の条件と ORDER BY の並び順を追えれば正解できる。
\end{advice}

% =============================================================
\section{正規化}
% =============================================================

\begin{teigi}[正規化の3段階]
\begin{itemize}
  \item \termtag{第1正規形}：繰り返し項目を排除（1つのセルに1つの値）
  \item \termtag{第2正規形}：部分関数従属を排除（主キーの一部だけに依存する列を別テーブルに分離）
  \item \termtag{第3正規形}：推移的関数従属を排除（主キー以外の列に依存する列を別テーブルに分離）
\end{itemize}
正規化の目的は\textbf{データの冗長性を排除}し、\textbf{更新時の矛盾を防ぐ}こと。
\end{teigi}

\begin{pitfall}[正規化のレベルを混同]
\begin{itemize}
  \item 第1：「繰り返しがない」だけではまだ第1
  \item 第2：「主キー全体に依存」＝部分関数従属がない
  \item 第3：「主キー以外への依存もない」＝推移的関数従属がない
  \item 「1→繰り返し除去、2→部分除去、3→推移除去」の語呂で覚える
\end{itemize}
\end{pitfall}

% =============================================================
\section{過去問ドリル}
% =============================================================

\shoumatsuhead{過去問ドリル：ネットワークとDB}

\begin{itpmon}[\sourcebadge{original}\enspace\diffbadge{標}]{%
次のSQL文を実行した結果として得られるものはどれか。

\texttt{SELECT 氏名, 点数 FROM 成績表 WHERE 点数 >= 80 ORDER BY 点数 DESC}}
\sentakushi
  {成績表から点数が80以上のレコードを点数の昇順で取得する。}
  {成績表から点数が80以上のレコードを点数の降順で取得する。}
  {成績表から点数が80未満のレコードを点数の昇順で取得する。}
  {成績表からすべてのレコードを点数の降順で取得する。}
\end{itpmon}
\begin{kaisetsu}{イ}
WHERE 点数 >= 80 で80点以上を絞り込み、ORDER BY 点数 \textbf{DESC}（降順）で並べ替え。DESCは大きい順、ASCは小さい順。
\end{kaisetsu}

\begin{itpmon}[\sourcebadge{original}\enspace\diffbadge{標}]{%
RDB（関係データベース）における主キーと外部キーの説明として、最も適切なものはどれか。}
\sentakushi
  {主キーはNULLを許容し、外部キーはNULLを許容しない。}
  {主キーはレコードを一意に識別し、外部キーは他テーブルの主キーを参照する。}
  {主キーも外部キーも、同一テーブル内でのみ使用される。}
  {外部キーは主キーの代わりに使うことができる。}
\end{itpmon}
\begin{kaisetsu}{イ}
\textbf{主キー}はレコードの一意識別（NULL・重複不可）、\textbf{外部キー}は他テーブルの主キーを参照してテーブル間を関連づける。アは逆。
\end{kaisetsu}

\begin{itpmon}[\sourcebadge{original}\enspace\diffbadge{強}]{%
サブネットマスクが255.255.255.0のネットワークで、ホストに割り当てられるIPアドレスの最大数はどれか。}
\sentakushi
  {252}
  {254}
  {255}
  {256}
\end{itpmon}
\begin{kaisetsu}{イ}
ホスト部は8ビットで $2^8 = 256$ 通りだが、ネットワークアドレス（全0）とブロードキャストアドレス（全1）を除くため $256 - 2 = \textbf{254}$。
\end{kaisetsu}

% =============================================================
\section{よくあるミスと到達確認}
% =============================================================

\begin{yokumiss}[この章でよくあるミス]
\begin{enumerate}[leftmargin=1.6em, itemsep=5pt]
  \item \textbf{TCPとUDPを逆にする}——TCP=信頼性重視、UDP=速度重視。「Tは丁寧（Teinei）」で覚える。
  \item \textbf{DNSとDHCPを混同する}——DNS=名前解決、DHCP=アドレス自動配布。
  \item \textbf{SQLのDESC/ASCを逆に覚える}——DESC=降順（大→小）、ASC=昇順（小→大）。
  \item \textbf{主キーにNULLが許される}と思い込む——主キーはNULL禁止・重複禁止。
  \item \textbf{サブネットマスクの計算で$-2$を忘れる}——ネットワークアドレスとブロードキャストは除外。
  \item \textbf{正規化のレベルを混同する}——1=繰り返し除去、2=部分従属除去、3=推移従属除去。
\end{enumerate}
\end{yokumiss}

\begin{tatsaku}
  \item TCP/IP 4層と各層のプロトコルを3つ以上言える
  \item TCPとUDPの違いを3点挙げられる
  \item DNS/DHCPの役割を即答できる
  \item サブネットマスクからホスト数を計算できる
  \item SQLのSELECT文を読んで結果を予測できる
  \item 主キーと外部キーの違いを説明できる
  \item 正規化の第1〜第3正規形を区別できる
  \item POP3とIMAPの違いを説明できる
\end{tatsaku}

\nextchapter{%
セキュリティとAIに進みます。CIA三要素、暗号方式、マルウェア対策、そしてAI/ML/DL——
テクノロジ系の最重要テーマを攻略します。
}
